\documentclass[UTF8,12pt]{ctexart}
\usepackage[paperwidth=240mm,paperheight=200mm,margin=14mm,top=8mm]{geometry}
\usepackage{amsmath,amssymb,mathtools}
\usepackage{xcolor}
\pagestyle{empty}
\setlength{\parindent}{0pt}
\setlength{\parskip}{0pt}
\newcommand{\qn}[1]{\textbf{\large #1}\ }
\xeCJKDeclareCharClass{CJK}{"2460 -> "24FF}
\begin{document}
\qn{1989.1}
$\lim\limits_{x\to0}x\cot 2x=$______．

\vfill
\newpage
\qn{1989.2}
$\int_0^\pi t\sin t\,\mathrm{d}t=$______．

\vfill
\newpage
\qn{1989.3}
曲线 $y=\int_0^x(t-1)(t-2)\,\mathrm{d}t$ 在点 $(0,0)$ 处的切线方程是______．

\vfill
\newpage
\qn{1989.4}
设 $f(x)=x(x+1)(x+2)\cdots(x+n)$，则 $f'(0)=$______．

\vfill
\newpage
\qn{1989.5}
设 $f(x)$ 是连续函数，且 $f(x)=x+2\int_0^1 f(t)\,\mathrm{d}t$，则 $f(x)=$______．

\vfill
\newpage
\qn{1989.6}
设 $f(x)=\begin{cases} a+bx^2, & x\leqslant 0, \\ \dfrac{\sin bx}{x}, & x>0 \end{cases}$ 在 $x=0$ 处连续，则常数 $a$ 与 $b$ 应满足的关系是______．

\vfill
\newpage
\qn{1989.7}
设 $\tan y=x+y$，则 $\mathrm{d}y=$______．

\vfill
\newpage
\qn{1989.8}
已知 $y=\arcsin\mathrm{e}^{-\sqrt{x}}$，求 $y'$．

\vfill
\newpage
\qn{1989.9}
求 $\int\dfrac{\mathrm{d}x}{x\ln^2x}$．

\vfill
\newpage
\qn{1989.10}
求 $\lim\limits_{x\to0}(2\sin x+\cos x)^{\frac{1}{x}}$．

\vfill
\newpage
\qn{1989.11}
已知 $\begin{cases} x=\ln(1+t^2), \\ y=\arctan t, \end{cases}$ 求 $\dfrac{\mathrm{d}y}{\mathrm{d}x}$，$\dfrac{\mathrm{d}^2y}{\mathrm{d}x^2}$．

\vfill
\newpage
\qn{1989.12}
已知 $f(2)=\dfrac{1}{2}$，$f'(2)=0$ 及 $\int_0^2 f(x)\,\mathrm{d}x=1$，求 $\int_0^1 x^2 f''(2x)\,\mathrm{d}x$．

\vfill
\newpage
\qn{1989.13}
当 $x>0$ 时，曲线 $y=x\sin\dfrac{1}{x}$（　　）
\\
\noindent (A) 有且仅有水平渐近线．　　(B) 有且仅有铅直渐近线．
\\
\noindent (C) 既有水平渐近线，也有铅直渐近线．　　(D) 既无水平渐近线，也无铅直渐近线．

\vfill
\newpage
\qn{1989.14}
若 $3a^2-5b<0$，则方程 $x^5+2ax^3+3bx+4c=0$（　　）
\\
\noindent (A) 无实根．　　(B) 有唯一实根．
\\
\noindent (C) 有三个不同实根．　　(D) 有五个不同实根．

\vfill
\newpage
\qn{1989.15}
曲线 $y=\cos x\left(-\dfrac{\pi}{2}\leqslant x\leqslant\dfrac{\pi}{2}\right)$ 与 $x$ 轴所围成的图形，绕 $x$ 轴旋转一周所成的旋转体的体积为（　　）
\\
\noindent (A) $\dfrac{\pi}{2}$．　　(B) $\pi$．
\\
\noindent (C) $\dfrac{\pi^2}{2}$．　　(D) $\pi^2$．

\vfill
\newpage
\qn{1989.16}
设两函数 $f(x)$ 和 $g(x)$ 都在 $x=a$ 处取得极大值，则函数 $F(x)=f(x)g(x)$ 在 $x=a$ 处（　　）
\\
\noindent (A) 必取极大值．　　(B) 必取极小值．
\\
\noindent (C) 不可能取极值．　　(D) 是否取极值不能确定．

\vfill
\newpage
\qn{1989.17}
微分方程 $y''-y=\mathrm{e}^x+1$ 的一个特解应具有形式（式中 $a,b$ 为常数）（　　）
\\
\noindent (A) $a\mathrm{e}^x+b$．　　(B) $ax\mathrm{e}^x+b$．
\\
\noindent (C) $a\mathrm{e}^x+bx$．　　(D) $ax\mathrm{e}^x+bx$．

\vfill
\newpage
\qn{1989.18}
设 $f(x)$ 在点 $x=a$ 的某个邻域内有定义，则 $f(x)$ 在 $x=a$ 处可导的一个充分条件是（　　）
\\
\noindent (A) $\lim\limits_{h\to+\infty}h\left[f\left(a+\dfrac{1}{h}\right)-f(a)\right]$ 存在．
\\
\noindent (B) $\lim\limits_{h\to0}\dfrac{f(a+2h)-f(a+h)}{h}$ 存在．
\\
\noindent (C) $\lim\limits_{h\to0}\dfrac{f(a+h)-f(a-h)}{2h}$ 存在．
\\
\noindent (D) $\lim\limits_{h\to0}\dfrac{f(a)-f(a-h)}{h}$ 存在．
求微分方程 $xy'+(1-x)y=\mathrm{e}^{2x}$（$0<x<+\infty$）满足 $y(1)=0$ 的特解．
设 $f(x)=\sin x-\int_0^x(x-t)f(t)\,\mathrm{d}t$，其中 $f$ 为连续函数，求 $f(x)$．
证明方程 $\ln x=\dfrac{x}{\mathrm{e}}-\int_0^\pi\sqrt{1-\cos 2x}\,\mathrm{d}x$ 在区间 $(0,+\infty)$ 内有且仅有两个不同实根．
对函数 $y=\dfrac{x+1}{x^2}$ 填写下表．
|  |  |
| --- | --- |
| 单调减少区间 |  |
| 单调增加区间 |  |
| 极值点 |  |
| 极值 |  |
| 凹区间 |  |
| 凸区间 |  |
| 拐点 |  |
| 渐近线 |  |
设抛物线 $y=ax^2+bx+c$ 过原点，当 $0\leqslant x\leqslant1$ 时 $y\geqslant0$，又已知该抛物线与 $x$ 轴及直线 $x=1$ 所围图形的面积为 $\dfrac{1}{3}$，试确定 $a,b,c$ 的值，使此图形绕 $x$ 轴旋转一周而成的旋转体的体积 $V$ 最小．

\vfill
\newpage
\end{document}